\chapter{Mathematical and Evaluation Framework}\label{ch:framework}

\section{Baseline nondimensionalization}
Let $L_0$ and $U_0$ denote reference length and velocity, and let $\rho_f$
and $\mu$ denote the fluid density and dynamic viscosity. The baseline
uses
\begin{equation}
 \vect{x}^{*}=\frac{\vect{x}}{L_0},\quad
 t^{*}=\frac{tU_0}{L_0},\quad
 \vect{u}^{*}=\frac{\vect{u}}{U_0},\quad
 p_v^{*}=\frac{pL_0}{\mu U_0},\quad
 \Rey=\frac{\rho_fU_0L_0}{\mu}.
 \label{eq:baseline_scales}
\end{equation}
Stars are suppressed after the scales have been specified. The steady
carrier equations are
\begin{equation}
 \nabla\cdot\vect{u}=0,\qquad
 \Rey(\vect{u}\cdot\nabla)\vect{u}
 =-\nabla p_v+\nabla^2\vect{u}.
 \label{eq:baseline_ns}
\end{equation}
Using a viscous pressure scale places $\Rey$ on the convective term.
Pressure values from this model cannot be compared directly with a model
using an inertial pressure scale without conversion.

For a spherical heavy particle, the retained force is the creeping-flow
drag. With $R=\rho_p/\rho_f$ and $d_p=d/L_0$, its dimensionless equations
are
\begin{equation}
 \frac{d\vect{q}}{dt}=\vect{v},\qquad
 \frac{d\vect{v}}{dt}
 =\frac{18}{R\Rey d_p^2}[\vect{u}(\vect{q})-\vect{v}]
 =\frac{\vect{u}(\vect{q})-\vect{v}}{\St},
 \qquad \St=\frac{R\Rey d_p^2}{18}.
 \label{eq:baseline_particle}
\end{equation}
The slip Reynolds number is
\begin{equation}
 \Rey_p=\Rey d_p|\vect{u}(\vect{q})-\vect{v}|.
 \label{eq:slip_re}
\end{equation}
The physical rationale for Stokes drag requires $\Rey_p\ll1$, along with a
small particle relative to the spatial scales of the carrier flow.

\section{Geometry and boundary-value matching}
The baseline channel is a constant-width tube around a planar centerline
$\vect{c}(s)$. Let $\vect{n}(s)$ denote its unit normal and let $h$ denote
the half-width. Coordinates are mapped by
\begin{equation}
 \vect{x}(s,\eta)=\vect{c}(s)+h\eta\vect{n}(s),\qquad
 0\leq s\leq L,\quad -1\leq\eta\leq1.
 \label{eq:centerline_map}
\end{equation}
The carrier velocity vanishes on both walls. At each end section, a
parabolic velocity in the local tangent direction is prescribed:
\begin{equation}
 \vect{u}(s_{\mathrm{end}},\eta)
 =(1-\eta^2)\vect{t}(s_{\mathrm{end}}).
 \label{eq:end_bc}
\end{equation}
Thus the reference problem uses velocity conditions at both ends, rather
than an imposed pressure outlet. A pressure gauge fixes the remaining
constant. The neural and finite-element comparisons must use this same
problem.

The associated volume flux per unit reference depth is
\begin{equation}
 Q_0=h\int_{-1}^{1}(1-\eta^2)\,d\eta=\frac{4h}{3}.
 \label{eq:flux}
\end{equation}
For a steady incompressible flow without sources of fluid mass, the section
flux should be independent of $s$. This makes flux variation a useful
integrated diagnostic alongside local differential residuals.

\section{Gaussian-channel scales}
For the Gaussian stage, the reference length is the channel height $H$.
The pressure scale is inertial:
\begin{equation}
 p_i^{*}=\frac{p}{\rho_f U_0^2},\qquad
 \Rey_H=\frac{U_0H}{\nu},\qquad \delta_p=\frac{d}{H}.
 \label{eq:gaussian_scales}
\end{equation}
For identical length and velocity scales, the two pressure conventions
satisfy $p_v^{*}=\Rey_H p_i^{*}$. The governing equation is correspondingly
written with $1/\Rey_H$ multiplying viscous diffusion.

The carrier is two-dimensional, but the represented physical particles are
spheres. A reference depth $B$ is therefore required to convert represented
sphere mass into a planar source. With $b=B/H$ and multiplicity $N_i$, the
dimensionless represented mass is
\begin{equation}
 m_i=\frac{N_iR\pi\delta_{p,i}^3}{6b},\qquad
 \St_i=\frac{R\Rey_H\delta_{p,i}^2}{18},\qquad
 \beta=\frac{\sum_i m_i}{|\Omega|}.
 \label{eq:parcel_mass}
\end{equation}
This convention scales physical mass by $\rho_fH^2B$. The kernel has units
of inverse area before nondimensionalization. Replacing it with a
three-dimensional volume kernel without changing the mass convention would
be dimensionally inconsistent.

\section{What is learned}
The neural representations approximate fields or trajectories; the
physical equations specify their constraints. The baseline force network
learns a known drag response from an algebraic residual. It does not infer
an unknown force law from experiments. The Gaussian parcel network predicts
time-dependent position and velocity corrections whose consistency is
enforced by both kinematic and dynamic residuals.

Adam is used for initial optimization, with optional limited-memory
quasi-Newton refinement \citep{kingma2015,liu1989}. The optimizer sequence
does not define the credibility of the final state. Checkpoint selection
uses a fixed residual monitor, whereas the reported field comparison uses
an independent reference. A longer training budget is a computational
choice, not an accuracy certificate.

\section{Error measures and pressure gauge}
For values at $M$ independent locations, the sample-relative velocity error
is
\begin{equation}
 E_u=\left[
 \frac{\sum_{j=1}^{M}|\vect{u}^{\mathrm{PINN}}_j-
               \vect{u}^{\mathrm{ref}}_j|^2}
      {\sum_{j=1}^{M}|\vect{u}^{\mathrm{ref}}_j|^2}
 \right]^{1/2}.
 \label{eq:velocity_error}
\end{equation}
These are discrete norms over the specified samples. Unless area weights
are explicitly included, they are not exact domain-integral norms. For
pressure, I first remove the arithmetic mean of each compared array:
\begin{equation}
 \widetilde p_j=p_j-\frac{1}{M}\sum_{k=1}^{M}p_k,\qquad
 E_p=\frac{\|\widetilde{\vect{p}}^{\mathrm{PINN}}-
                    \widetilde{\vect{p}}^{\mathrm{ref}}\|_2}
                   {\|\widetilde{\vect{p}}^{\mathrm{ref}}\|_2}.
 \label{eq:pressure_error}
\end{equation}
The pressure samples span the streamwise direction so that the reference
denominator is meaningful. Mean subtraction removes only a constant; it
does not rescale a pressure gradient or modify the predicted dynamics.

The section-flux diagnostic is
\begin{equation}
 E_Q=\max_s\left|\frac{Q(s)}{Q_0}-1\right|.
 \label{eq:flux_error}
\end{equation}
Reported differential residuals use their stated nondimensional equations.
Their magnitudes should not be compared across pressure scales, equation
normalizations, or geometries without accounting for those differences.

\section{Trajectory evaluation}
Consider a neural and reference trajectory with the same initial condition.
If their common interval is $[0,T_c]$, the position discrepancy is measured
over that interval and normalized by channel width $W=2h$. Endpoint
agreement is assessed separately using the terminal event class. I also
retain the fraction of the longer trajectory covered by the common
interval. A low error calculated over a short common interval is not
sufficient evidence of agreement over the full transit.

Integration success has a narrower meaning: it indicates that the numerical
ODE solver completed its assigned calculation without reporting failure.
It does not establish that the neural velocity or force was accurate.
Accordingly, the acceptance criteria distinguish integration status,
position discrepancy, common-time coverage, and endpoint agreement.

\section{Evidence hierarchy}
I use four levels of numerical evidence. Algebraic identities establish
properties of the coupling operators. Analytical benchmarks test reduced
equations with known solutions. Independently discretized references test
the field approximation for a matched problem. Finally, transport
comparisons assess the accumulated consequences of the field and force
errors. None of these alone establishes experimental agreement.

For the baseline, the numerical screening limits are 5\% for field and
normalized drag errors, 2\% for boundary RMS and reference-interpolation
error, and 5\% for maximum flux deviation. Baseline particle acceptance is
stricter in its combination of conditions: every tested endpoint class must
agree, every width-normalized RMS position error must be below 5\%, and
every common-time fraction must exceed 95\%. For the Gaussian verification
tests, the stated relative-error and independent residual criteria are
0.03. These limits are study criteria rather than universal physical
thresholds. Their actual outcomes are reported where the corresponding
results are presented.
